\section{Introdução}


Em 1750, B. Franklin sugeriu a existência de partículas elétricas. 
Mais de 100 anos depois, em 1859, J Plücker descobriu os raios catódicos enquanto realizava experimentos 
sobre condução elétrica em gases. 
Seu aluno, J. Hittorf, descobriu em 1869 que a trajetória desses raios era afetada pela presença de campos magnéticos
assim como são linhas de corrente elétrica. \cite{cathoderays}

Realizando experimentos com os raios catódicos, J. Thomson inferiu que átomos sao compostos de partículas subatômicas 
com massa 1000 vezes menor que a massa do átomo, e com carga elétrica negativa. 
Essas particulas posteriormente receberam o nome de elétron. 
Em 1897, J. Thomson e P. Zeeman independentemente determinaram experimentalmente a razão carga-massa 
($e/m$) do elétron. \cite{nobelmillikan}

O próximo passo nessa jornada foi, naturalmente, determinar experimentalmente a carga e/ou a massa do elétron. 
Essa missão foi aceita por R. Millikan e H. Fletcher em 1909, e os resultados foram publicados em 1913. 
Eles encontraram que a carga do elétron era $\qty{-1,592e-19}{\coulomb}$, o que corresponde a um desvio de $0,6\%$ 
com relação ao valor aceito atualmente de $\qty{-1,602e-19}{\coulomb}$. \cite{millikan}

Nesse trabalho, é utilizado um aparato similar ao utilizado por Millikan e Fletcher em 1909 para 
determinar o valor da carga fundamental.